Each \(B_c^{(m)}\) is a nonzero scalar multiple of a monomial of total degree \(m\), so \(V_m\) consists of restrictions of polynomials of degree at most \(m\).  Conversely, on the simplex \(\sum_{a,y}q_{ay}=1\).  For a monomial \(q^{\gamma}\) of degree \(d\leq m\),
\[
 q^{\gamma}=q^{\gamma}\left(\sum_{a,y}q_{ay}\right)^{m-d},
\]
and the multinomial expansion is a linear combination of total-degree-\(m\) monomials, hence of the Bernstein polynomials.  Thus \(V_m\) is exactly the asserted restriction space.

Suppose a polynomial \(P\) agreed with \(f\) on \(K_{\epsilon}\).  Put \(a=\epsilon/2>0\), and for \(r\in(\epsilon,1-\epsilon)\) define
\[
 q(r)=(1-r,0,r-a,a).
\]
All coordinates are nonnegative, they sum to one, and \(p(q(r))=r\), so the path lies in \(K_{\epsilon}\).  The function \(R(r)=P(q(r))\) is a polynomial, while
\[
 f(q(r))=\frac ar.
\]
Thus the polynomial \(rR(r)-a\) vanishes on a nonempty open interval and is identically zero.  Evaluation at \(r=0\) gives \(-a=0\), a contradiction.  Hence \(f\) is not the restriction of any polynomial and, in particular, \(f\notin V_m\) for every finite \(m\).